\documentclass[aspectratio=169,12pt]{beamer}
\usepackage{amsmath}
\usepackage{amssymb}
\usepackage{array}
\usepackage[most]{tcolorbox}
\usepackage{graphicx}
\usepackage[sfdefault,lf]{FiraSans}
\renewcommand{\familydefault}{\sfdefault}
\setlength{\parindent}{0pt}
\everymath{\displaystyle}

\setbeamertemplate{navigation symbols}{}
\setbeamertemplate{footline}{}
\setbeamertemplate{headline}{}
\setbeamertemplate{blocks}[rounded][shadow=false]

\definecolor{mmpageWhite}{HTML}{FCFBF7}
\definecolor{mmtanSoft}{HTML}{F7F5EF}
\definecolor{mmgreenDeep}{HTML}{244B2C}
\definecolor{mmgreenLeaf}{HTML}{5D8A56}
\definecolor{mmgreenSoft}{HTML}{E4EFD9}
\definecolor{mmtext}{HTML}{163221}
\definecolor{mmwarn}{HTML}{8F2F36}
\definecolor{mmwarnSoft}{HTML}{F8E8EA}

\setbeamercolor{background canvas}{bg=mmpageWhite}
\setbeamercolor{normal text}{fg=mmtext,bg=mmpageWhite}
\setbeamercolor{block title}{fg=mmgreenDeep,bg=mmgreenSoft}
\setbeamercolor{block body}{fg=mmtext,bg=mmtanSoft}
\setbeamercolor{block title example}{fg=mmgreenDeep,bg=mmgreenSoft}
\setbeamercolor{block body example}{fg=mmtext,bg=white}
\setbeamercolor{block title alerted}{fg=mmwarn,bg=mmwarnSoft}
\setbeamercolor{block body alerted}{fg=mmtext,bg=white}
\setbeamercolor{itemize item}{fg=mmgreenDeep}
\setbeamercolor{itemize subitem}{fg=mmgreenDeep}

\newtcolorbox{MMCard}[2][]{enhanced,breakable=false,colback=#2,colframe=black,boxrule=1.5pt,arc=8pt,left=5pt,right=5pt,top=4pt,bottom=4pt,before skip=0pt,after skip=0pt,#1}
\newcommand{\KoruIcon}{\raisebox{-0.2em}{\includegraphics[height=1.05em]{../../../manamaths/SITE/header-logo.png}}}
\newcommand{\NotesTitle}[1]{{\colorbox{mmtanSoft}{\parbox{0.975\linewidth}{\vspace{0.14em}\hspace{0.25em}{\LARGE\bfseries\textcolor{mmgreenDeep}{#1}\hfill\KoruIcon}\vspace{0.14em}}}\par\vspace{0.18em}{\color{mmgreenLeaf}\rule{\linewidth}{2.0pt}}\par\vspace{0.12em}}}
\newcommand{\SectionTitle}[1]{{\large\bfseries\textcolor{mmgreenDeep}{#1}}\par\vspace{0.04em}}
\newcommand{\GreenDivider}{\par\vspace{0.01em}{\color{mmgreenLeaf}\rule{\linewidth}{2.4pt}}\par\vspace{0.03em}}
\newcommand{\KeyIdea}[1]{\begin{MMCard}[title={\textbf{Key idea}},colbacktitle=mmgreenSoft,coltitle=mmgreenDeep]{mmtanSoft}\small #1\end{MMCard}}
\newcommand{\WorkedExample}[2]{\begin{MMCard}[title={\textbf{#1}},colbacktitle=mmgreenSoft,coltitle=mmgreenDeep]{white}\small #2\end{MMCard}}
\newcommand{\CommonMistake}[1]{\begin{MMCard}[title={\textbf{Common mistake}},colbacktitle=mmwarnSoft,coltitle=mmwarn]{white}\small #1\end{MMCard}}
\newcommand{\TryThese}[1]{\begin{MMCard}[title={\textbf{Try these}},colbacktitle=mmgreenSoft,coltitle=mmgreenDeep]{mmtanSoft}\small #1\end{MMCard}}
\newcommand{\StepLine}[2]{\textbf{#1.} #2\par\vspace{0.03em}}

\usepackage{tikz}

\setbeamertemplate{background}{
 \begin{tikzpicture}[remember picture,overlay]
 \draw[
 black,
 line width=1.5pt,
 rounded corners=10pt
 ]
 ([xshift=0.45cm,yshift=-0.45cm]current page.north west)
 rectangle
 ([xshift=-0.45cm,yshift=0.45cm]current page.south east);
 \end{tikzpicture}
}

\begin{document}

\begin{frame}[t,shrink=8]
\NotesTitle{Notes \& Steps}
\footnotesize
\begin{columns}[T,onlytextwidth]
\begin{column}{0.48\textwidth}
\KeyIdea{Integers are whole numbers that can be positive, negative, or zero. They have no fractions or decimals. Use a number line to visualise addition and subtraction.}
\SectionTitle{Steps — adding integers}
\StepLine{1}{Start at the first number on the number line.}
\StepLine{2}{If adding a positive, move right. If adding a negative, move left.}
\SectionTitle{Steps — subtracting integers}
\StepLine{1}{Change subtraction to adding the opposite.}
\StepLine{2}{Then follow the addition rules.}
\StepLine{3}{Example: \(7 - (-2) = 7 + 2 = 9\).}
\end{column}
\begin{column}{0.48\textwidth}
\SectionTitle{Key facts}
\begin{itemize}
  \setlength{\itemsep}{0.12em}
  \item Opposite of \(-7\) is \(7\).
  \item \(3, -5, 0\) are integers; \(1/2\) and \(-5.5\) are not.
  \item \(-3 + 5 = 2\) (start at \(-3\), move 5 right)
  \item \(7 - (-2) = 7 + 2 = 9\)
  \item \((-3) + (-4) = -7\)
  \item On a number line, \(-1 > -5\) because \(-1\) is further right.
\end{itemize}
\end{column}
\end{columns}
\GreenDivider
\vfill
\begin{columns}[b,onlytextwidth]
\begin{column}{0.48\textwidth}
\CommonMistake{Thinking \(-5\) is greater than \(-1\). On the number line, \(-1\) is to the right of \(-5\), so \(-1 > -5\).}
\end{column}
\begin{column}{0.48\textwidth}
\TryThese{\begin{enumerate}
  \item What is the opposite of \(-7\)?
  \item Calculate: \(-3 + 5\).
  \item Order from smallest: \(-2, 0, 3\).
\end{enumerate}}
\end{column}
\end{columns}
\end{frame}

\begin{frame}[t,shrink=12]
\NotesTitle{Notes \& Steps}
\begin{columns}[T,onlytextwidth]
\begin{column}{0.48\textwidth}
\WorkedExample{Example 1: temperature rise}{A temperature is \(-4^\circ\)C. It rises by \(10^\circ\)C. What is the new temperature?\[-4 + 10 = 6\]Answer: \(6^\circ\)C.}
\end{column}
\begin{column}{0.48\textwidth}
\WorkedExample{Example 2: adding negatives}{Calculate \((-3) + (-4)\).\[(-3) + (-4)\]Start at \(-3\), move 4 left: \(-7\). Answer: \(-7\).}
\end{column}
\end{columns}
\GreenDivider
\TryThese{\begin{enumerate}
  \item Calculate: \(5 + (-2)\).
  \item Calculate: \(-1 - 3\).
  \item Calculate: \(7 - (-2)\).
\end{enumerate}}
\end{frame}

\begin{frame}[t,shrink=12]
\NotesTitle{Notes \& Steps}
\begin{columns}[T,onlytextwidth]
\begin{column}{0.48\textwidth}
\WorkedExample{Example 3: multiply negatives}{Calculate \((-3) \times 4\).\[(-3) \times 4\]Positive \(\times\) negative = negative. Answer: \(-12\).}
\end{column}
\begin{column}{0.48\textwidth}
\WorkedExample{Example 4: diving scenario}{A diver at \(-12\) m ascends 3 m then descends 5 m.\[-12 + 3 - 5 = -14\]Answer: final depth is \(-14\) m.}
\end{column}
\end{columns}
\GreenDivider
\CommonMistake{Forgetting that two negatives make a positive. \((-3) \times (-4) = 12\), not \(-12\). Think: subtracting a debt means gaining money.}
\end{frame}

\end{document}
